(1) 定义 \(\Omega_R=\sqrt{\Omega^2+\Delta^2}\) 和单位矢量 \(\hat n=(\Omega,0,\Delta)/\Omega_R\)。因 \((\hat n\cdot\boldsymbol\sigma)^2=I\)，指数演化为 \(U(t)=\cos(\Omega_Rt/2)I-i\sin(\Omega_Rt/2)(\Delta\sigma_z+\Omega\sigma_x)/\Omega_R\)。
(2) 由上态到下态的振幅是 \(\langle\downarrow|U(t)|\uparrow\rangle=-i(\Omega/\Omega_R)\sin(\Omega_Rt/2)\)，因此 \(P_\downarrow(t)=\Omega^2/(\Omega^2+\Delta^2)\sin^2(\Omega_Rt/2)\)。广义Rabi频率 \(\Omega_R\) 同时决定振荡周期与失谐导致的振荡对比度。
(3) 最大翻转概率为 \(P_{\max}=\Omega^2/(\Omega^2+\Delta^2)\)，只有 \(\Delta=0\) 时才能达到一。共振时 \(P_\downarrow=\sin^2(\Omega t/2)\)，首次完全翻转发生在 \(t_\pi=\pi/\Omega\)，称为 \(\pi\) 脉冲；非零失谐时纵向哈密顿量分量使演化轴倾斜，无法从北极精确旋转至南极。
(4) 旋转参考系哈密顿量不随时间改变，故 \(\langle H\rangle\) 守恒，初态期望为 \(\hbar\Delta/2\)。在实验室系中驱动电磁场可与物质交换能量，旋转波近似产生的有效守恒准能不应被误称为外驱动完全不做功。
